% Charakterisierung der durch Leerwort und Adjunktion erzeugten Menge.
% Vorausgesetzt: generated-set sowie die früheren Mengen-, Graphen- und Folgenresultate.
% Keine spätere Wortaxiomatik.
\subsection{Die erzeugten Mengen sind genau die endlichen Folgen}

Die Definition der Wortmenge verwendet nur das Leerwort und einen
Erzeugungsschritt. Jetzt zeigen wir, welche Mengen dabei entstehen:
Genau die endlichen Folgen aus Band~21 gehören zur erzeugten Menge.
Die lückenlose Reihenfolge ist somit eine Folgerung aus den Regeln.

\Needspace{16\baselineskip}
\noindent\begin{minipage}{\linewidth}
\FormulaThmDeltaK[Die Erzeugungsregeln liefern lückenlose Funktionsgraphen]{%
  w\in\FiniteWords A\vdash w\colon\NatSeg{\FiniteCard(w)}\to A
}{GeneratedWordGraphTyping}{%
  \DeltaRow{Mengen}{A\dsep w\dsep u\dsep a}
  \DeltaRow{Abkürzung}{P(u)\coloneqq
    u\colon\NatSeg{\FiniteCard(u)}\to A}}
\end{minipage}\par
Für Wortargumente \(u\in\FiniteWords A\) ist die Endlichkeit bereits
durch \FormulaRefAuto{FiniteWordUnderlyingFinite}[thm] gesichert.
Die Kardinalzahl in \(P(u)\) ist damit definiert. Im Anfangsfall verwenden
wir die bereits bestimmte Kardinalzahl der leeren Menge.
\Needspace{11\baselineskip}
\begin{tabproofsplitwide}
  \proofpartwideindR[Anfang der Erzeugung]{P(\varnothing)}{%
    \varnothing\colon\NatSeg{\FiniteCard(\varnothing)}\to A
  }[GeneratedWordGraphTypingBase]
    \proofstepwidestar[]{\varnothing\colon\varnothing\to A}{%
      \FormulaRefAuto{\varnothing\colon \varnothing\to M}[thm]}
    \proofstepwidestar[]{\NatSeg0=\varnothing}{\FormulaRefAuto{PeanoNatSegZero}[thm]}
    \proofstepwidestar[]{\FiniteCard(\varnothing)=0}{\FormulaRefAuto{B27FiniteCardEmpty}[thm]}
    \proofstepwidestar[]{P(\varnothing)}{\rIE{2,3,1}}
  \closeproofpartwideindR
  \Needspace{11\baselineskip}
  \proofpartwideindR[Ein Erzeugungsschritt]{%
    u\in\FiniteWords A\dsep a\in A\dsep P(u)\vdash P(\OccurrenceAdjunction u a)}{%
    \begin{aligned}[t]
      &u\in\FiniteWords A\dsep a\in A\dsep
        u\colon\NatSeg{\FiniteCard(u)}\to A\vdash{}\\[-2pt]
      &\OccurrenceAdjunction u a\colon\NatSeg{\FiniteCard(\OccurrenceAdjunction u a)}\to A
    \end{aligned}
  }[GeneratedWordGraphTypingStep]
    \proofstepwidestar[1]{u\in\FiniteWords A}{\rA}
    \proofstepwidestar[2]{a\in A}{\rA}
    \proofstepwidestar[3]{P(u)}{\rA}
    \proofstepwidestar[1]{\Finite u}{\FormulaRefAuto{FiniteWordUnderlyingFinite}[thm]{1}}
    \proofstepwidestar[1]{\FiniteCard(u)\in\mathbb N}{\FormulaRefAuto{FiniteCardNatural}[thm]{4}}
    \proofstepwidestar[1,2,3]{u\cup\{(\FiniteCard(u),a)\}\colon\NatSeg{\FiniteCard(u)+1}\to A}{%
      \FormulaRefAuto{EarlyWordAdjunctionGraphTyping}[thm]{5,3,2}}
    \proofstepwidestar[1]{u\in\FiniteOccurrenceSets A}{\FormulaRefAuto{FiniteWordInUniverse}[thm]{1}}
    \proofstepwidestar[1,2]{\OccurrenceAdjunction u a=u\cup\{(\FiniteCard(u),a)\}}{%
      \FormulaRefAuto{OccurrenceCardinalAdjunctionDef}[def]{7,2}}
    \proofstepwidestar[1,2,3]{\OccurrenceAdjunction u a\colon\NatSeg{\FiniteCard(u)+1}\to A}{\rIE{8,6}}
    \proofstepwidestar[1]{\FiniteCard(u)+1\in\mathbb N}{\FormulaRefAuto{PeanoAddOneClosure}[thm]{5}}
    \proofstepwidestar[1,2,3]{\FiniteCard(\OccurrenceAdjunction u a)=\FiniteCard(u)+1}{%
      \FormulaRefAuto{FiniteWordCardinalityFromTyping}[thm]{10,9}}
    \proofstepwidestar[1,2,3]{P(\OccurrenceAdjunction u a)}{\rIE{11,9}}
  \closeproofpartwideindR
  \Needspace{7\baselineskip}
  \proofpartwide{Induktionsschluss}
    \proofstepwidestar[1]{w\in\FiniteWords A}{\rA}
    \proofstepwidestar[1]{w\colon\NatSeg{\FiniteCard(w)}\to A}{\FormulaRefAuto{GeneratedWordInductionRule}[thm]{IA,IS,1}}
  \closeproofpartwide
\end{tabproofsplitwide}

\noindent\begin{minipage}{\linewidth}
\FormulaThmDeltaK[Jede endliche Folge wird durch die Wortregeln erzeugt]{%
  n\in\mathbb N\dsep w\colon\NatSeg n\to A\vdash w\in\FiniteWords A
}{B27FiniteGraphGenerated}{%
  \DeltaRow{Mengen}{A\dsep n\dsep k\dsep w\dsep u}
  \DeltaRow{Abkürzungen}{Q(k)\coloneqq
    \forall u\,(u\colon\NatSeg k\to A\rightarrow u\in\FiniteWords A)
    \dsep r\coloneqq u\restriction_{\NatSeg k}}}
\end{minipage}\par
Hier wird die Zahleninduktion aus Band~10 benutzt: Ein Graph über
\(\NatSeg{k+1}\) entsteht aus seiner Einschränkung auf \(\NatSeg k\)
durch genau den zugelassenen Erzeugungsschritt.
\Needspace{11\baselineskip}
\begin{tabproofsplitwide}
  \proofpartwideindR[Induktionsanfang]{Q(0)}{%
    \forall u\,(u\colon\NatSeg0\to A\rightarrow u\in\FiniteWords A)
  }[B27FiniteGraphGeneratedBase]
    \proofstepwidestar[1]{u\colon\NatSeg0\to A}{\rA}
    \proofstepwidestar[]{\NatSeg0=\varnothing}{\FormulaRefAuto{PeanoNatSegZero}[thm]}
    \proofstepwidestar[1]{u\colon\varnothing\to A}{\rIE{2,1}}
    \proofstepwidestar[1]{u=\varnothing}{\FormulaRefAuto{f\colon \varnothing\to M \vdash f=\varnothing}[thm]{3}}
    \proofstepwidestar[]{\varnothing\in\FiniteWords A}{\FormulaRefAuto{GeneratedWordEmpty}[thm]}
    \proofstepwidestar[1]{u\in\FiniteWords A}{\rIE{4,5}}
    \proofstepwidestar[]{Q(0)}{\rUI{\rRI{1,6}}}
  \closeproofpartwideindR
  \Needspace{11\baselineskip}
  \proofpartwideindR[Induktionsschritt]{k\in\mathbb N\dsep Q(k)\vdash Q(k+1)}{%
    \begin{aligned}[t]
      &k\in\mathbb N\dsep
        \forall u\,(u\colon\NatSeg k\to A\rightarrow u\in\FiniteWords A)\vdash{}\\[-2pt]
      &\forall u\,(u\colon\NatSeg{k+1}\to A\rightarrow u\in\FiniteWords A)
    \end{aligned}
  }[B27FiniteGraphGeneratedStep]
    \proofstepwidestar[1]{k\in\mathbb N}{\rA}
    \proofstepwidestar[2]{Q(k)}{\rA}
    \proofstepwidestar[3]{u\colon\NatSeg{k+1}\to A}{\rA}
    \proofstepwidestar[1]{\NatSeg k\subseteq\NatSeg{k+1}}{\FormulaRefAuto{PeanoNatSegSubsetAddOne}[thm]{1}}
    \proofstepwidestar[1,3]{r\colon\NatSeg k\to A}{\FormulaRefAuto{RestrictionDef}[def]{3,4}}
    \proofstepwidestar[1,2,3]{r\in\FiniteWords A}{\rRE{5,\rUE{2}}}
    \proofstepwidestar[1,3]{\FiniteCard(r)=k}{\FormulaRefAuto{FiniteWordCardinalityFromTyping}[thm]{1,5}}
    \proofstepwidestar[1]{k\in\NatSeg{k+1}}{\FormulaRefAuto{PeanoNatSegEndpointInAddOne}[thm]{1}}
    \proofstepwidestar[1,3]{u(k)\in A}{\FormulaRefAuto{F\colon A\to B\dsep x\in A\vdash F(x)\in B}[thm]{3,8}}
    \proofstepwidestar[1,2,3]{\OccurrenceAdjunction r{u(k)}\in\FiniteWords A}{\FormulaRefAuto{GeneratedWordClosed}[thm]{6,9}}
    \proofstepwidestar[1,2,3]{r\in\FiniteOccurrenceSets A}{\FormulaRefAuto{FiniteWordInUniverse}[thm]{6}}
    \proofstepwidestar[1,2,3]{\OccurrenceAdjunction r{u(k)}=r\cup\{(\FiniteCard(r),u(k))\}}{\FormulaRefAuto{OccurrenceCardinalAdjunctionDef}[def]{11,9}}
    \proofstepwidestar[1,3]{u=r\cup\{(k,u(k))\}}{\FormulaRefAuto{B27FiniteGraphLastAdjunction}[thm]{1,3}}
    \proofstepwidestar[1,2,3]{u=\OccurrenceAdjunction r{u(k)}}{\rIE{7,12,13}}
    \proofstepwidestar[1,2,3]{u\in\FiniteWords A}{\rIE{14,10}}
    \proofstepwidestar[1,2]{Q(k+1)}{\rUI{\rRI{3,15}}}
  \closeproofpartwideindR
  \Needspace{7\baselineskip}
  \proofpartwide{Schluss durch Zahleninduktion}
    \proofstepwidestar[1]{n\in\mathbb N}{\rA}
    \proofstepwidestar[2]{w\colon\NatSeg n\to A}{\rA}
    \proofstepwidestar[1]{Q(n)}{\FormulaRefAuto{PeanoInductionRuleAddOne}[thm]{IA,IS,1}}
    \proofstepwidestar[1,2]{w\in\FiniteWords A}{\rRE{2,\rUE{3}}}
  \closeproofpartwide
\end{tabproofsplitwide}

\Needspace{17\baselineskip}
\noindent\begin{minipage}{\linewidth}
\FormulaThmDeltaK[Funktionenkriterium der endlichen Wortmenge]{%
  w\in\FiniteWords A\eqvdash
  \Bigl(w\in\powerset(\mathbb N\times A)\land
    \exists n\in\mathbb N\,(w\colon\NatSeg n\to A)\Bigr)
}{FiniteWordSetMembership}{\DeltaRow{Mengen}{A\dsep w\dsep n}}
\end{minipage}\par
\Needspace{7\baselineskip}
\begin{tabproofsplitwide}
  \proofpartwide{Von der Erzeugung zum endlichen Funktionsgraphen}
    \proofstepwidestar[1]{w\in\FiniteWords A}{\rA}
    \proofstepwidestar[1]{w\in\FiniteOccurrenceSets A}{\FormulaRefAuto{FiniteWordInUniverse}[thm]{1}}
    \proofstepwidestar[1]{w\subseteq\mathbb N\times A\land\Finite w}{\FormulaRefAuto{FiniteOccurrenceSetUniverseMembership}[thm]{2}}
    \proofstepwidestar[1]{w\in\powerset(\mathbb N\times A)}{\FormulaRefAuto{x\in\powerset(A)\eqvdash x\subseteq A}[thm]{\rAEa{3}}}
    \proofstepwidestar[1]{\FiniteCard(w)\in\mathbb N}{\FormulaRefAuto{FiniteCardNatural}[thm]{\rAEb{3}}}
    \proofstepwidestar[1]{w\colon\NatSeg{\FiniteCard(w)}\to A}{\FormulaRefAuto{GeneratedWordGraphTyping}[thm]{1}}
    \proofstepwidestar[1]{\exists n\in\mathbb N\,(w\colon\NatSeg n\to A)}{\rEI{\rAI{5,6}}}
    \proofstepwidestar[1]{w\in\powerset(\mathbb N\times A)\land
      \exists n\in\mathbb N\,(w\colon\NatSeg n\to A)}{\rAI{4,7}}
  \closeproofpartwide
  \Needspace{7\baselineskip}
  \proofpartwide{Vom endlichen Funktionsgraphen zur Erzeugung}
    \proofstepwidestar[1]{w\in\powerset(\mathbb N\times A)\land
      \exists n\in\mathbb N\,(w\colon\NatSeg n\to A)}{\rA}
    \proofstepwidestar[2]{n\in\mathbb N\land w\colon\NatSeg n\to A}{\rA}
    \proofstepwidestar[2]{w\in\FiniteWords A}{\FormulaRefAuto{B27FiniteGraphGenerated}[thm]{\rAEa{2},\rAEb{2}}}
    \proofstepwidestar[1]{w\in\FiniteWords A}{\rEE{\rAEb{1},2,3}}
  \closeproofpartwide
\end{tabproofsplitwide}

Damit darf ein Wort jetzt als endliche Folge gelesen werden. Die
Folgenbeschreibung ersetzt die Erzeugungsdefinition nicht, sondern
beschreibt genau ihre Elemente und stellt die Positionsnotation bereit.


